\begin{enumerate}
    \item Définissez la notion de radiation monochromatique.
    \item Dans le tableau ci-dessous, on a indiqué la fréquence de quelques
          radiations lumineuses monochromatiques. Calculez leur longueur d'onde
          dans le vide et précisez leur couleur ou leur domaine d'appartenance~:

          \begin{minipage}{\linewidth}
              \begin{table}[H]
              \centering
              \setlength{\arrayrulewidth}{0.8pt}
              %\setlength{\tabcolsep}{18pt}
              \renewcommand{\arraystretch}{1.5}
              \begin{tabularx}{\textwidth}{|C|C|C|C|C|C|C|}
                  \hline
                  $\nu~(\unit{\Hz})$ & $\num{5.45e14}$ & $\num{4.00e14}$ &
                  $\num{8.36e14}$ & $\num{3.42e14}$ & $\num{7.14e14}$ &
                  $\num{5.00e14}$ \\
                  \hline
                  $\lambda~(\unit{\nm})$ &  &  &  &  &  &  \\
                  \hline
                  \begin{tabular}{c}
                      Couleur ou \\[-1mm]
                      domaine \\
                  \end{tabular} &  &  &  &  &  &  \\
                  \hline
              \end{tabularx}
          \end{table}
          \end{minipage}

          \textbf{Rq}~: On pourrait exprimer ces très grandes fréquences grâce à
          un multiple du Hertz, le terahertz~:
          $\qty{1}{\THz}=\qty{e12}{\Hz}$
\end{enumerate}

\begin{donnees}
    Célérité de la lumière dans le vide~: $c=\qty{300000}{\km\per\s}$
\end{donnees}

